%----------------------------------------------------------------------------------------
% STEEKPROEVEN EN SAMENHANG. Alle opgaven zijn eigen werk.
%
% Opbouw: Start, leerstof (theorie, werkwijze of uitleg, voorbeeld, korte oefening), en op een
% nieuwe pagina de Opdrachten (door elkaar, stijgende moeilijkheid), Koppel en Test jezelf.
% Labels van dit hoofdstuk beginnen met h1-.
%----------------------------------------------------------------------------------------

\mychapter{1}{Steekproeven en samenhang}

\lkr{Leerplan B+S'': LPD 50 (randomisatie en representativiteit bij steekproeven, steekproefvariabiliteit, aselecte steekproef) en LPD 51 (samenhang tegenover causaliteit), met de wenken over selectiebias, andere steekproefmethodes, simulaties, observationeel en experimenteel onderzoek. Spreidingsdiagrammen en de correlatiecoëfficiënt kennen leerlingen uit de 2de graad; hier gaat het om de interpretatie. Dit hoofdstuk heeft weinig rekenwerk: laat de leerlingen vooral argumenteren, in volzinnen.}

Hoeveel jongeren lezen elke dag? Werkt een nieuw medicijn? Gaan leerlingen die meer slapen, beter presteren? Je kunt niet iedereen ondervragen of testen. Daarom onderzoek je een steekproef, en trek je daaruit een besluit over de hele groep. Dat lukt alleen als de steekproef goed gekozen is, en als je voorzichtig bent met wat je besluit. In dit hoofdstuk leer je hoe je een steekproef kiest, waarom resultaten van steekproef tot steekproef verschillen, en waarom samenhang nog geen oorzaak is.

\begin{oefeningen}[Start]
\oef{h1-start-1}

\oef{h1-verken}
\end{oefeningen}

\lkr{Laat oefening \ref{oef:h1-verken} maken vóór de uitleg, het liefst met de hele klas: elke leerling of elk groepje trekt een eigen steekproef met het rekentoestel. Zet de gemiddelden op het bord. Zo zien ze steekproefvariabiliteit voor ze het woord kennen.}

%========================================================================================
\section{Populatie en steekproef}

\begin{theorie}
De \textbf{\textcolor{accent}{populatie}} is de hele groep waarover je iets wilt weten. Een \textbf{\textcolor{accent}{steekproef}} is een deel van de populatie dat je echt onderzoekt.

\vspace{1mm}
Een getal dat de populatie beschrijft, zoals het populatiegemiddelde $\mu$ of de populatieproportie $p$, heet een \textbf{\textcolor{accent}{parameter}}. Het overeenkomstige getal uit de steekproef, zoals $\formule{\xbar}$ of $\formule{\hat{p}}$, is een \textbf{\textcolor{accent}{steekproefgrootheid}}.

\vspace{1mm}
Een steekproef is \textbf{\textcolor{accent}{representatief}} als ze een goede afspiegeling is van de populatie. Bij een \textbf{\textcolor{accent}{aselecte}} steekproef heeft elk lid van de populatie dezelfde kans om gekozen te worden.
\end{theorie}

Je werkt met een steekproef omdat het onmogelijk is iedereen te bereiken, omdat het te lang duurt of te veel kost, of omdat het onderzoek de onderzochte dingen vernielt, zoals bij het testen van de levensduur van lampen. Een steekproefgrootheid gebruik je als schatting van de parameter: $\xbar$ schat $\mu$, en $\hat{p}$ (\textit{p hoed}) schat $p$.

Representativiteit kun je niet rechtstreeks garanderen: je kent de populatie niet volledig. Je bereikt ze door bewust het toeval te laten kiezen. Dat heet \textbf{randomisatie}: lootjes trekken, of de leden nummeren en willekeurige nummers laten kiezen door een rekentoestel.

\voorbeeld{}{Een school wil weten hoeveel uur haar 600 leerlingen per week sporten. Ze laat de computer 60 leerlingnummers kiezen, en die leerlingen vullen een vragenlijst in. Gemiddeld sporten ze 3,4 uur. Wat is de populatie, de steekproef, de parameter en de steekproefgrootheid?\label{vb:h1-begrippen}}{%
\uitleg{We zoeken welke groep de school wil kennen, welke groep ze echt onderzoekt, en welk getal bij elk hoort.}
\begin{stappenbreed}
\text{populatie: de 600 leerlingen} & De groep waarover de school iets wil weten. \\[.6ex]
\text{steekproef: de 60 gekozen leerlingen} & Aselect: elk nummer had dezelfde kans. \\[.6ex]
\text{parameter } \mu \text{: gemiddeld aantal uren van alle 600} & Onbekend. \\[.6ex]
\text{steekproefgrootheid } \xbar = 3{,}4 & Een schatting van $\mu$.
\end{stappenbreed}}

\begin{oefeningen}
\oef{h1-begrippen}
\end{oefeningen}

%========================================================================================
\section{Vertekende steekproeven}

Een steekproef die niet representatief is, geeft een vertekend beeld, hoe groot ze ook is. Die fout heet \textbf{bias} of \textbf{vertekening}. De bekendste oorzaken:

{\renewcommand{\arraystretch}{1.4}%
\begin{tabular}{@{}>{\RaggedRight\arraybackslash\bfseries}p{3.6cm}>{\RaggedRight\arraybackslash}p{12cm}@{}}
selectiebias & Sommige leden hebben meer kans om in de steekproef te komen: een enquête aan de schoolpoort om 8 uur mist wie met de bus later komt. \\
vrijwilligers & Wie zelf kiest om mee te doen, heeft vaak een uitgesproken mening: een online peiling op een sportsite over sport. \\
non-respons & Wie niet antwoordt, kan anders zijn dan wie wel antwoordt: drukke mensen beantwoorden minder enquêtes over vrije tijd. \\
vraagstelling & Een sturende vraag stuurt het antwoord: \textit{Vind jij ook dat de school te weinig doet voor het klimaat?}
\end{tabular}\par}

Een grote steekproef is niet vanzelf representatief. Een online peiling met 100\,000 stemmen kan sterk vertekend zijn, terwijl een aselecte steekproef van 1000 personen een goede schatting geeft. De grootte helpt pas als de steekproef aselect is.

\voorbeeld{}{Een gemeente wil weten of haar inwoners een nieuw zwembad willen. Ze stelt de vraag aan de bezoekers van het huidige zwembad. 85~\% is voor. Is dat een goede schatting? Hoe zou je het beter doen?\label{vb:h1-bias}}{%
\uitleg{We vragen ons af of elke inwoner dezelfde kans had om ondervraagd te worden, en of de ondervraagden anders kunnen zijn dan de anderen.}
De steekproef is niet representatief: zwembadbezoekers zwemmen graag, en zijn dus waarschijnlijk vaker voor een nieuw zwembad dan de gemiddelde inwoner (selectiebias). De 85~\% overschat het aandeel voorstanders. Beter: kies aselect inwoners uit het bevolkingsregister, en vraag het hun op een neutrale manier.}

Naast de enkelvoudige aselecte steekproef zijn er andere methodes, die ook het toeval gebruiken:

\begin{itemize}
\item een \textbf{gestratificeerde} steekproef: je verdeelt de populatie in groepen (strata), zoals de drie graden van een school, en neemt uit elke groep een aselecte steekproef, in verhouding tot de grootte van de groep;
\item een \textbf{getrapte} steekproef: je kiest eerst aselect enkele groepen, zoals scholen, en daarbinnen aselect personen.
\end{itemize}

\voorbeeld{}{De school uit voorbeeld \ref{vb:h1-begrippen} telt 200 leerlingen in de eerste graad, 220 in de tweede en 180 in de derde. Hoe kies je een gestratificeerde steekproef van 60 leerlingen?\label{vb:h1-strata}}{%
\uitleg{Elke graad krijgt evenveel plaatsen in verhouding tot zijn aantal leerlingen: telkens 60 op 600, dus een tiende.}
\begin{stappenbreed}
200 \cdot \tfrac{60}{600} = 20, \quad 220 \cdot \tfrac{60}{600} = 22, \quad 180 \cdot \tfrac{60}{600} = 18 & Per graad een tiende. \\[1ex]
\text{binnen elke graad aselect kiezen} & Bijvoorbeeld met \texttt{randInt} op het rekentoestel.
\end{stappenbreed}
Zo is elke graad zeker vertegenwoordigd in de juiste verhouding.}

\begin{oefeningen}
\oef{h1-bias}
\end{oefeningen}

%========================================================================================
\section{Steekproefvariabiliteit}

In de Verken gaven drie aselecte steekproeven uit dezelfde klas drie verschillende gemiddelden. Dat is geen fout: het is \textbf{steekproefvariabiliteit}. Een steekproefgrootheid hangt af van het toeval, en is dus een toevalsvariabele.

\naast{\begin{minipage}{7.4cm}
De figuur toont een simulatie. In een populatie heeft 40~\% een bepaalde eigenschap. De computer trok 40 aselecte steekproeven van 20 personen, en 40 van 200 personen, en berekende telkens het aandeel $\hat{p}$. Elke stip is een steekproef.

Bij $n = 20$ liggen de aandelen tussen 0,10 en 0,60; bij $n = 200$ tussen ongeveer 0,30 en 0,52. Grotere aselecte steekproeven schommelen minder rond de parameter $p = 0{,}4$. In hoofdstuk 3 en 4 zie je hoe groot die schommeling is.
\end{minipage}}{\begin{tikzpicture}[x=11cm,y=.16cm,font=\scriptsize]
  \begin{scope}[yshift=2.6cm]
  \draw[grijs] (.05,0) -- (.65,0);
  \foreach \p in {(0.1,1),(0.15,1),(0.2,1),(0.2,2),(0.2,3),(0.25,1),(0.25,2),(0.25,3),(0.3,1),(0.3,2),(0.3,3),(0.35,1),(0.35,2),(0.35,3),(0.35,4),(0.35,5),(0.35,6),(0.35,7),(0.35,8),(0.35,9),(0.35,10),(0.4,1),(0.4,2),(0.4,3),(0.4,4),(0.4,5),(0.45,1),(0.45,2),(0.45,3),(0.45,4),(0.5,1),(0.5,2),(0.5,3),(0.5,4),(0.5,5),(0.55,1),(0.55,2),(0.55,3),(0.55,4),(0.6,1)}{\fill[accent] \p circle (1.5pt);}
  \node[anchor=west] at (.05,11.5) {$n = 20$};
  \end{scope}
  \draw[grijs] (.05,0) -- (.65,0);
  \foreach \p in {(0.3,1),(0.325,1),(0.35,1),(0.35,2),(0.35,3),(0.35,4),(0.375,1),(0.375,2),(0.375,3),(0.375,4),(0.375,5),(0.375,6),(0.375,7),(0.375,8),(0.4,1),(0.4,2),(0.4,3),(0.4,4),(0.4,5),(0.4,6),(0.4,7),(0.4,8),(0.4,9),(0.4,10),(0.4,11),(0.425,1),(0.425,2),(0.425,3),(0.425,4),(0.425,5),(0.425,6),(0.425,7),(0.45,1),(0.45,2),(0.45,3),(0.45,4),(0.475,1),(0.475,2),(0.475,3),(0.525,1)}{\fill[fel!80!black] \p circle (1.5pt);}
  \node[anchor=west] at (.05,11.5) {$n = 200$};
  \foreach \t/\l in {.1/{0,1},.2/{0,2},.3/{0,3},.4/{0,4},.5/{0,5},.6/{0,6}}{\draw[grijs] (\t,0) -- (\t,-.5) node[below,text=black] {\l};}
  \node[below] at (.35,-3.5) {aandeel $\hat{p}$ in de steekproef};
\end{tikzpicture}}

%========================================================================================
\section{Samenhang en causaliteit}

\subsection{Samenhang is geen oorzaak}

Uit een spreidingsdiagram zie je of twee variabelen samenhangen: stijgt de ene meestal als de andere stijgt, dan is er een positieve samenhang. De correlatiecoëfficiënt $r$, tussen $-1$ en $1$, meet hoe sterk een lineaire samenhang is. Maar samenhang zegt niet dat de ene variabele de andere veroorzaakt.

\begin{theorie}
Twee variabelen hangen \textbf{\textcolor{accent}{samen}} (zijn \textbf{\textcolor{accent}{gecorreleerd}}) als hun waarden samen variëren. Er is \textbf{\textcolor{accent}{causaliteit}} als een verandering van de ene variabele een verandering van de andere veroorzaakt.

\vspace{1mm}
Samenhang zonder causaliteit kan ontstaan door

\vspace{1mm}
\quad 1.\ \ een \textbf{\textcolor{accent}{verborgen variabele}} die beide beïnvloedt;

\quad 2.\ \ \textbf{\textcolor{accent}{omgekeerde causaliteit}}: het gevolg is eigenlijk de oorzaak;

\quad 3.\ \ \textbf{\textcolor{accent}{toeval}}, zeker bij kleine aantallen of als je veel variabelen vergelijkt.
\end{theorie}

\voorbeeld{}{Over een jaar hangen de verkoop van zonnecrème en het aantal huisartsbezoeken voor hooikoorts sterk samen. Veroorzaakt zonnecrème hooikoorts?\label{vb:h1-verborgen}}{%
\uitleg{We zoeken een derde variabele die beide kan verklaren.}
Nee. In de lente en de zomer schijnt de zon meer, waardoor mensen meer zonnecrème kopen, en is er meer pollen in de lucht, waardoor er meer hooikoorts is. Het seizoen is een verborgen variabele: het verklaart de samenhang, zonder dat het ene het andere veroorzaakt.}

\voorbeeld{}{Uit een onderzoek blijkt dat mensen die veel naar de huisarts gaan, minder gezond zijn dan wie weinig gaat. Besluit iemand: \textit{huisartsbezoeken maken je ziek}. Wat is er fout?\label{vb:h1-omgekeerd}}{%
\uitleg{We kijken of oorzaak en gevolg niet omgekeerd kunnen zijn.}
Wie ziek is, gaat vaker naar de huisarts: de slechte gezondheid is de oorzaak van de bezoeken, niet het gevolg. Dat is omgekeerde causaliteit.}

\subsection{Observationeel en experimenteel onderzoek}

Bij \textbf{observationeel onderzoek} kijk je alleen: je vraagt mensen hoeveel ze slapen en welke punten ze halen. Een samenhang kan dan altijd door een verborgen variabele komen. Bij \textbf{experimenteel onderzoek} grijp je zelf in: je verdeelt de deelnemers \textit{aselect} in een testgroep, die een behandeling krijgt, en een controlegroep, die ze niet krijgt (of een nepbehandeling, een placebo). Door de aselecte verdeling zijn de groepen gemiddeld gelijk in alle andere opzichten. Een verschil tussen de groepen kun je dan aan de behandeling toeschrijven.

\voorbeeld{}{Hoe zou je onderzoeken of naar muziek luisteren tijdens het studeren de toetsresultaten verbetert?\label{vb:h1-experiment}}{%
\uitleg{Alleen vragen wie met muziek studeert, is observationeel: wie met muziek studeert, kan in andere opzichten verschillen. We ontwerpen een experiment.}
Neem een groep vrijwillige leerlingen en verdeel ze door loting in twee groepen. Beide groepen studeren dezelfde tekst, even lang, in hetzelfde lokaal: de ene groep met muziek, de andere zonder. Daarna maken ze dezelfde toets. Een duidelijk verschil in de gemiddelde score kun je aan de muziek toeschrijven, omdat de loting de andere verschillen gemiddeld wegwerkt. Of het verschil \textit{duidelijk} genoeg is, leer je in hoofdstuk 3.}

Een experiment is niet altijd mogelijk of toegelaten. Of roken schadelijk is, kun je niet testen door mensen te laten roken; dan blijft observationeel onderzoek over, met veel voorzichtigheid en met veel verschillende studies.

\begin{werkwijze}
Zie je een samenhang, vraag je dan af: kan een derde variabele beide verklaren? Kan het omgekeerd zijn? Kan het toeval zijn? Was het onderzoek observationeel of een experiment met een aselecte verdeling in groepen? Alleen bij een goed experiment mag je over oorzaak spreken.
\end{werkwijze}

\begin{oefeningen}
\oef{h1-causaliteit}
\end{oefeningen}

%========================================================================================
\opdrachtenkop

\begin{oefeningen}
\oef{h1-o1}

\oef{h1-o2}

\oef{h1-o3}

\oef{h1-o4}

\oef{h1-o5}

\oef{h1-o6}

\oef{h1-o7}

\oef{h1-o8}

\oef{h1-o9}

\oef{h1-o10}
\end{oefeningen}

%----------------------------------------------------------------------------------------
% KOPPEL: verbanden met andere leerstof
%----------------------------------------------------------------------------------------
\needspace{20\baselineskip}
\opdrachtsectie{Koppel}

\lkr{De eerste oefening gebruikt de verwachtingswaarde en de standaardafwijking van een binomiale verdeling (Kansrekenen, zelfde jaar), de tweede de regressierechte uit de 2de graad.}

\begin{oefeningen}
\oef{h1-koppel-1}

\oef{h1-koppel-2}
\end{oefeningen}

%----------------------------------------------------------------------------------------
\opdrachtsectie{Test jezelf}

\begin{oefeningen}
\oef{h1-t1}

\oef{h1-t2}

\oef{h1-t3}

\oef{h1-t4}
\end{oefeningen}

\herhaalschema{%
  \fv{oef:h1-t1}{vb:h1-begrippen}{oef:h1-begrippen}%
  \fv{oef:h1-t2}{vb:h1-bias}{oef:h1-bias}%
  \fv{oef:h1-t3}{vb:h1-strata}{oef:h1-bias}%
  \fv{oef:h1-t4}{vb:h1-verborgen}{oef:h1-causaliteit}}

\ifzelfstudie
\section*{Maak je eigen samenvatting}
Sluit dit boek en schrijf in je schrift: wat een populatie, een steekproef, een parameter en een steekproefgrootheid zijn, wat representatief en aselect betekenen, de oorzaken van bias, gestratificeerde en getrapte steekproeven, wat steekproefvariabiliteit is en hoe ze afhangt van de grootte, en de drie oorzaken van samenhang zonder causaliteit, met het verschil tussen observationeel en experimenteel onderzoek. Vergelijk daarna met de kaders.
\fi
